\begin{abstract}
In prediction markets, are historical winners merely fortunate, or do they possess superior forecasting ability? Can wallet trading behavior provide additional predictive information beyond market implied probabilities? This study investigates Smart Money in Polymarket by examining the persistence of forecasting skills, the predictive value of trading behavior, and changes in betting size.

Using approximately 570 million historical trades from January 2025 to March 2026, this study constructs a dataset covering 1.17 million wallets and 14,925 markets. Permutation tests and Empirical Bayes estimation are employed to examine forecasting skills, while Logistic Regression, LightGBM, MLP, Mean Pooling, and GRU models are developed to evaluate the predictive value of wallet behavior.

The results reveal statistically significant persistence in forecasting performance ($p\approx 0.007$), with stronger signals observed among experienced wallets. Empirical Bayes estimation improves the future forecasting performance of the selected top 10\% wallets by approximately 44\%. Further analysis shows that incorporating behavioral and historical trading features improves probability forecasts, while controlled experiments provide encouraging evidence of the additional information contained in historical trading behavior.

An intriguing finding concerns changes in betting size. Wallets tend to increase their stakes on outcomes already favored by the market. However, the largest additional forecasting advantage associated with larger bets occurs when market probabilities are between 40\% and 60\%. Within this interval, wallets placing unusually large bets outperform those maintaining similar stakes by 3.05 and 3.69 percentage points across two evaluation periods. These findings suggest that identifying Smart Money requires consideration of both historical forecasting performance and the circumstances under which wallets adjust their trading behavior.
\end{abstract}
